
\documentclass{article}
\usepackage{colortbl}
\usepackage{makecell}
\usepackage{multirow}
\usepackage{supertabular}

\begin{document}

\newcounter{utterance}

\centering \large Interaction Transcript for game `codenames', experiment `association\_difficult', episode 6 with reflexion{-}llama{-}3.1{-}8b{-}instruct.
\vspace{24pt}

{ \footnotesize  \setcounter{utterance}{1}
\setlength{\tabcolsep}{0pt}
\begin{supertabular}{c@{$\;$}|p{.15\linewidth}@{}p{.15\linewidth}p{.15\linewidth}p{.15\linewidth}p{.15\linewidth}p{.15\linewidth}}
    \# & $\;$A & \multicolumn{4}{c}{Game Master} & $\;\:$B\\
    \hline

    \theutterance \stepcounter{utterance}  
    & & & \multicolumn{2}{p{0.3\linewidth}}{
        \cellcolor[rgb]{0.9,0.9,0.9}{
            \makecell[{{p{\linewidth}}}]{
                \texttt{\tiny{[GM$|$GM]}}
                {'hidden': {'team': ['mint', 'masher', 'daffodil', 'aunt'], 'innocent': ['tulip', 'peony', 'millimeter', 'ostrich', 'father'], 'opponent': [], 'assassin': ['bowling']}, 'revealed': {'team': {'team': [], 'innocent': [], 'opponent': [], 'assassin': []}, 'opponent': {'team': [], 'innocent': [], 'opponent': [], 'assassin': []}}}
            }
        }
    }
    & & \\ \\

    \theutterance \stepcounter{utterance}  
    & & \multicolumn{4}{p{0.6\linewidth}}{
        \cellcolor[rgb]{0.9,0.9,0.9}{
            \makecell[{{p{\linewidth}}}]{
                \texttt{\tiny{[P1$\langle$GM]}}
                \texttt{You have attempted this same task before and did not succeed. These are the plans you wrote after those attempts.} \\
\texttt{Trial \#1: New plan: } \\
\\ 
\texttt{In my previous attempts, I was able to guess some of the team's words, but I failed to guess all of them before the opposing team uncovered all of theirs. Upon reviewing my attempts, I noticed that I was focusing on individual words and their associations, but I didn't consider the overall strategy of the game. } \\
\\ 
\texttt{In the future, I will take a more holistic approach by considering the relationships between the words and the clues. I will also pay closer attention to the words that have been removed from the lists and try to find connections between the remaining words and the clues. Additionally, I will make sure to provide more general and abstract clues that can relate to multiple words, rather than focusing on specific words or their literal meanings. } \\
\\ 
\texttt{For example, in my next attempt, I might provide a clue that relates to the theme of the words, such as "season" for words like "turkey" and "daffodil", or "tool" for words like "masher" and "bicycle". This will allow me to cover more ground and increase the chances of guessing all of the team's words before the opposing team does.} \\
\\ 
\texttt{Trial \#2: New plan: } \\
\\ 
\texttt{In my previous attempts, I was able to provide a clue that related to some of the team's words, but I failed to guess all of them before the opposing team uncovered all of theirs. Upon reviewing my attempts, I noticed that I was focusing on individual words and their associations, but I didn't consider the overall strategy of the game. } \\
\\ 
\texttt{In the future, I will take a more holistic approach by considering the relationships between the words and the clues. I will also pay closer attention to the words that have been removed from the lists and try to find connections between the remaining words and the clues. Additionally, I will make sure to provide more general and abstract clues that can relate to multiple words, rather than focusing on specific words or their literal meanings. } \\
\\ 
\texttt{To achieve this, I will start by analyzing the team's words and identifying common themes, categories, or relationships between them. I will then use this information to create a clue that is broad enough to relate to multiple words, but specific enough to be helpful. } \\
\\ 
\texttt{For example, in my next attempt, I might provide a clue that relates to the theme of the words, such as "garden" for words like "daffodil" and "mint", or "tool" for words like "masher" and "bicycle". This will allow me to cover more ground and increase the chances of guessing all of the team's words before the opposing team does.} \\
\\ 
\\ 
\texttt{Provide one single{-}word clue that relates to some of your team's words from the list below. You can choose to how many and to which words your clue relates to.} \\
\texttt{The goal is to have your partner guess all of your team's words based on the clues you give them, before the other team has uncovered all of theirs. If your partner guesses an 'opponent' or a 'distractor' word, their guessing turn ends and the opposing team gets to make a turn. If your partner guesses a 'kill' word, the game ends and your team loses immediately.} \\
\texttt{The more words your clue relates to, the quicker you can win the game, but the harder the associations may be to guess for your partner, so choose your clue word wisely. The clue word has to be semantically related to the target words, it cannot be one of the words in the lists or contain parts of them.} \\
\\ 
\texttt{Always give your single{-}word clue and your comma{-}separated list of related target words in the following format and make your answers as short as possible, never include any other text than is required in this form:} \\
\\ 
\texttt{CLUE: <WORD>} \\
\texttt{TARGETS: <TARGETS>} \\
\\ 
\texttt{Your team words are: mint, masher, daffodil, aunt.} \\
\texttt{Your opponent's team words are: .} \\
\texttt{Distractor words are: tulip, peony, millimeter, ostrich, father.} \\
\texttt{The kill words are: bowling.} \\
            }
        }
    }
    & \\ \\

    \theutterance \stepcounter{utterance}  
    & \multicolumn{4}{p{0.6\linewidth}}{
        \cellcolor[rgb]{0.9,0.9,0.9}{
            \makecell[{{p{\linewidth}}}]{
                \texttt{\tiny{[P1$\rangle$GM]}}
                \texttt{CLUE: garden} \\
\texttt{TARGETS: daffodil, mint, masher} \\
            }
        }
    }
    & & \\ \\

    \theutterance \stepcounter{utterance}  
    & & & \multicolumn{2}{p{0.3\linewidth}}{
        \cellcolor[rgb]{0.9,0.9,0.9}{
            \makecell[{{p{\linewidth}}}]{
                \texttt{\tiny{[GM$|$GM]}}
                \texttt{garden} \\
            }
        }
    }
    & & \\ \\

    \theutterance \stepcounter{utterance}  
    & & & \multicolumn{2}{p{0.3\linewidth}}{
        \cellcolor[rgb]{0.9,0.9,0.9}{
            \makecell[{{p{\linewidth}}}]{
                \texttt{\tiny{[GM$|$GM]}}
                ['daffodil', 'mint', 'masher']
            }
        }
    }
    & & \\ \\

    \theutterance \stepcounter{utterance}  
    & & & \multicolumn{2}{p{0.3\linewidth}}{
        \cellcolor[rgb]{0.9,0.9,0.9}{
            \makecell[{{p{\linewidth}}}]{
                \texttt{\tiny{[GM$|$GM]}}
                {'word': 'daffodil', 'assignment': 'team'}
            }
        }
    }
    & & \\ \\

    \theutterance \stepcounter{utterance}  
    & & & \multicolumn{2}{p{0.3\linewidth}}{
        \cellcolor[rgb]{0.9,0.9,0.9}{
            \makecell[{{p{\linewidth}}}]{
                \texttt{\tiny{[GM$|$GM]}}
                {'word': 'mint', 'assignment': 'team'}
            }
        }
    }
    & & \\ \\

    \theutterance \stepcounter{utterance}  
    & & & \multicolumn{2}{p{0.3\linewidth}}{
        \cellcolor[rgb]{0.9,0.9,0.9}{
            \makecell[{{p{\linewidth}}}]{
                \texttt{\tiny{[GM$|$GM]}}
                {'word': 'masher', 'assignment': 'team'}
            }
        }
    }
    & & \\ \\

    \theutterance \stepcounter{utterance}  
    & & & \multicolumn{2}{p{0.3\linewidth}}{
        \cellcolor[rgb]{0.9,0.9,0.9}{
            \makecell[{{p{\linewidth}}}]{
                \texttt{\tiny{[GM$|$GM]}}
                \texttt{opponent has won} \\
            }
        }
    }
    & & \\ \\

    \theutterance \stepcounter{utterance}  
    & \multicolumn{4}{p{0.6\linewidth}}{
        \cellcolor[rgb]{0.95,0.95,0.95}{
            \makecell[{{p{\linewidth}}}]{
                \texttt{\tiny{[P2$\rightarrow$P1]}}
                \texttt{Post{-}episode Reflexion (not sent to GM):} \\
\texttt{New plan: } \\
\\ 
\texttt{In my previous attempts, I was able to provide a clue that related to some of the team's words, but I failed to guess all of them before the opposing team uncovered all of theirs. Upon reviewing my attempts, I noticed that I was focusing on individual words and their associations, but I didn't consider the overall strategy of the game. } \\
\\ 
\texttt{In the future, I will take a more holistic approach by considering the relationships between the words and the clues. I will also pay closer attention to the words that have been removed from the lists and try to find connections between the remaining words and the clues. Additionally, I will make sure to provide more general and abstract clues that can relate to multiple words, rather than focusing on specific words or their literal meanings. } \\
\\ 
\texttt{To achieve this, I will start by analyzing the team's words and identifying common themes, categories, or relationships between them. I will then use this information to create a clue that is broad enough to relate to multiple words, but specific enough to be helpful. } \\
\\ 
\texttt{For example, in my next attempt, I might provide a clue that relates to the theme of the words, such as "garden" for words like "daffodil" and "mint", or "tool" for words like "masher" and "bicycle". This will allow me to cover more ground and increase the chances of guessing all of the team's words before the opposing team does.} \\
\\ 
\texttt{However, I also noticed that I provided a clue that was too specific and related to only a few words. In the future, I will make sure to provide a clue that is more general and can relate to multiple words. For example, I might provide a clue like "plant" or "object" that can relate to multiple words on the list.} \\
\\ 
\texttt{I will also pay closer attention to the distractor words and try to find connections between them and the team's words. This will help me to avoid providing clues that are too similar to the distractor words and increase the chances of guessing all of the team's words before the opposing team does.} \\
            }
        }
    }
    & & \\ \\

\end{supertabular}
}

\end{document}
